\chapter{Loss Optimization under the Throughput Constraint}

\section{Optimization Problem}

Once throughput is treated as intrinsic to the detector objective, the next
question is selection rather than discovery: among throughput-preserving
objectives, which
one gives the best detector utility? The \texttt{0403-combined-6strat-pitchrescale}
campaign was designed with six candidate losses, although only five appear in
the completed sweep summary.

\section{Cross-Strategy Comparison}

\runfigure{0403-combined-6strat-pitchrescale-cn2-5e14/19_cross_strategy_received_power.png}
  {0403 cross-strategy received power}
  {Mean detector-side received power at 10\,$\mu$m across the completed 0403
  strategies.}

\runfigure{0403-combined-6strat-pitchrescale-cn2-5e14/20_encircled_energy_per_strategy.png}
  {0403 encircled energy}
  {Encircled-energy comparison for the same strategy family.}

\begin{table}[H]
  \centering
  \begin{tabular}{lccc}
    \toprule
    \smalltablehead{Strategy} & \smalltablehead{$\PIB@10\,\mu$m} & \smalltablehead{$\TP$} & \smalltablehead{$\RP@10\,\mu$m} \\
    \midrule
    Turbulent baseline & 77.11\% & 97.00\% & 212.04 \\
    AbsBucket+CO & 80.52\% & 98.30\% & 224.28 \\
    PIB+HardTP(10) & 80.42\% & 98.92\% & 224.97 \\
    TP-PIB $w=0.5$ & 81.32\% & 96.94\% & 223.78 \\
    TP-PIB $w=2.0$ & 80.41\% & 98.93\% & 224.96 \\
    Raw RP & 80.83\% & 98.59\% & 225.54 \\
    \bottomrule
  \end{tabular}
  \caption{Completed 0403 strategy summary from
  \texttt{19\_cross\_strategy\_summary.json}. The designed sixth strategy
  \texttt{TP-PIB w=5.0} does not appear in the completed sweep summary.}
\end{table}

\section{Why the Direct Received-Power Objective Wins}

The winning objective is the direct received-power loss,
\begin{equation}
L_{\mathrm{RP}} = -\log\!\left(\frac{E_{\mathrm{bucket}}}{E_{\mathrm{input}}} + \epsilon\right).
\end{equation}
Its advantage is threefold:
\begin{itemize}
  \item it optimizes the detector quantity directly rather than a proxy ratio;
  \item the logarithm amplifies gradients in the low-RP regime, making early
  training less inert;
  \item throughput is embedded automatically because energy scattered out of the
  bucket immediately reduces the numerator.
\end{itemize}

Strong throughput penalties can approximate the same behavior, but the direct
received-power objective removes one layer of proxy design. That is why it
edges out the near-tie alternatives in mean detector performance.

\runfigure{0403-combined-6strat-pitchrescale-cn2-5e14/focal_raw_received_power/18_received_power_histogram.png}
  {Direct received-power histogram}
  {Detector-side received-power histogram for the mean-winning direct
  objective.}

\runfigure{0403-combined-6strat-pitchrescale-cn2-5e14/pib_hard_tp/18_received_power_histogram.png}
  {Hard-TP histogram}
  {A near-tie strategy: PIB with a strong throughput penalty.}

\runfigure{0403-combined-6strat-pitchrescale-cn2-5e14/focal_tp_pib_w2/18_received_power_histogram.png}
  {TP-PIB w=2 histogram}
  {Another near-tie strategy with strong stability and high throughput.}

\section{Interpreting Near-Ties}

Although \texttt{focal\_raw\_received\_power} has the highest mean received power, the best
stability-oriented interpretation is subtler. \texttt{PIB+HardTP(10)} and
\texttt{TP-PIB w=2.0} are essentially tied in mean detector performance and can be read as
more conservative alternatives. This matters because the objective family is no
longer about ``whether'' throughput matters, but about how strongly it should be
regularized.

\runfigure{0403-combined-6strat-pitchrescale-cn2-5e14/focal_raw_received_power/06_fig1_focal_plane_vacuum_vs_turbulent_vs_d2nn.png}
  {Direct received-power focal plane}
  {Focal-plane comparison for the final mean-winning direct received-power
  strategy.}

\runfigure{0403-combined-6strat-pitchrescale-cn2-5e14/focal_raw_received_power/14_zernike_delta_per_mode.png}
  {Raw-RP Zernike deltas}
  {Output-plane diagnostics included only to show that detector gains do not
  imply full wavefront recovery.}

\section{Handoff to the Distance Sweep}

This chapter closes the detector-objective design loop. The next experiment
freezes the direct received-power objective and asks a different question: how does the same
detector-centric optical strategy behave as turbulence strength increases with
range?
